\documentclass[11pt]{article}

\usepackage[margin=1in]{geometry}
\usepackage{amsmath,amssymb,amsthm,array,mathtools,microtype}
\usepackage[T1]{fontenc}
\usepackage[hidelinks]{hyperref}

\newtheorem{theorem}{Theorem}[section]
\newtheorem{proposition}[theorem]{Proposition}
\newtheorem{corollary}[theorem]{Corollary}
\newtheorem{lemma}[theorem]{Lemma}
\newcommand{\A}{\mathbb A}
\newcommand{\Spec}{\operatorname{Spec}}
\newcommand{\gdeg}{\operatorname{gdeg}}
\numberwithin{equation}{section}

\title{The Exact Nonproperness Locus of the Quadratic-Gauge Keller Maps}
\author{Roy van Rijn\\{\small Independent researcher}}
\date{27 July 2026}
\hypersetup{
  pdftitle={The Exact Nonproperness Locus of the Quadratic-Gauge Keller Maps},
  pdfauthor={Roy van Rijn},
  pdfsubject={Nonproperness and boundary sheets of polynomial Keller maps},
  pdfkeywords={Keller map, nonproperness locus, discriminant, Newton polygon,
    Jelonek locus}
}

\begin{document}
\maketitle

\begin{abstract}
For the all-degree quadratic-gauge Keller maps, we identify the reduced
nonproperness locus with the discriminant of the scalar inverse equation.
We compute the complete fiber table on the degenerate slice \(\Pi=0\) and
the exact global \(\Pi\)-adic order of the discriminant.  These geometric
results are independent of the prescribed finite-\'etale-fiber theorem.
\end{abstract}

\section{The quadratic-gauge input}

Let \(K\) be a characteristic-zero field and
\begin{equation}
 G(S)=g_1S+g_2S^2+\cdots+g_NS^N,\qquad g_1g_3g_N\neq0.
 \label{eq:seed}
\end{equation}
The quadratic-gauge map \(F_G=(\Pi,B,C):\A^3_K\to\A^3_K\) is the map
constructed in the companion prescribed-fiber paper.  Its scalar inverse
equation is
\[
 E_{\Pi,B,C}(S)=G_\Pi(S)-\frac{g_1}{2}(BS^2+C),
\]
where
\[
 G_\Pi(S)=g_1S+\Pi(g_2S^2+g_3S^3)
          +\sum_{k=4}^Ng_k\Pi^kS^k.
\]
It satisfies
\[
 \det DF_G=-2,\qquad \gdeg(F_G)=N.
\]
For \(\pi\neq0\), its literal fiber is
\begin{equation}
 F_G^{-1}(\pi,b,c)\simeq
 \Spec\left((K[S]/(E_{\pi,b,c}))[1/\overline{E_{\pi,b,c}'}]\right).
 \label{eq:localized-input}
\end{equation}
These identities follow from the compact inverse design and
root-to-source reconstruction in the prescribed-fiber paper; the present
note begins with them and studies the boundary.

For a squarefree \(P\in K[T]\) of degree \(N\geq3\) and
\(a\in K\) with \(P'(a)P'''(a)\neq0\), the prescribed-fiber seed is
\[
 G(S)=P(a+S)-P(a).
\]
After the determinant-one output normalization, the distinguished target is
\begin{equation}
 y_{P,a}=\left(1,0,-\frac{2P(a)}{P'(a)}\right),
 \label{eq:prescribed-target}
\end{equation}
and its inverse equation is \(P(a+S)\).

\input{sections/nonproperness}

\section*{Verification status}

The graph-boundary base-change argument, Jelonek descent, complete
\(\Pi=0\) fiber table, and Newton-polygon discriminant calculation are
ordinary mathematical proofs and are not part of the Lean formalization of
the prescribed-fiber theorem.  The exact symbolic audit is
\[
 \texttt{scripts/verify\_quadratic\_gauge\_nonproperness.py}.
\]

\begin{thebibliography}{9}
\footnotesize

\bibitem{Jelonek1993}
Z.~Jelonek,
\emph{The set of points at which a polynomial map is not proper},
Ann. Polon. Math. \textbf{58} (1993), 259--266,
\href{https://doi.org/10.4064/ap-58-3-259-266}
{doi:10.4064/ap-58-3-259-266}.

\end{thebibliography}

\end{document}
